% =====================================================================
%  Pacotes e configuracoes gerais do preambulo.
%  Adicione aqui os pacotes que o seu artigo precisar.
% =====================================================================

% --- Codificacao e idioma --------------------------------------------
\usepackage[T1]{fontenc}      % fontes com acentuacao correta e hifenizacao
\usepackage[utf8]{inputenc}   % o arquivo-fonte esta em UTF-8
\usepackage{lmodern}          % Latin Modern, versao vetorial da Computer Modern

% --- Elementos graficos ----------------------------------------------
\usepackage{graphicx}         % \includegraphics
\graphicspath{{figures/}}     % imagens sao procuradas em figures/
\usepackage{float}            % posicionamento [H]
\usepackage{booktabs}         % \toprule, \midrule, \bottomrule
\usepackage{multirow}
\usepackage{longtable}        % tabelas que quebram entre paginas

% --- Matematica ------------------------------------------------------
\usepackage{amsmath}
\usepackage{amssymb}

% --- Tipografia ------------------------------------------------------
\usepackage{microtype}        % ajuste fino de espacamento e protrusao
\usepackage{indentfirst}      % recua tambem o primeiro paragrafo da secao
\usepackage{nomencl}          % lista de simbolos
\usepackage{color}
\usepackage{xcolor}

% --- Links -----------------------------------------------------------
% A abntex2 ja carrega o hyperref; aqui apenas ajustamos a aparencia.
\definecolor{lgLink}{RGB}{0,80,160}
\hypersetup{
  colorlinks=true,
  linkcolor=lgLink,
  citecolor=lgLink,
  urlcolor=lgLink,
  bookmarksdepth=4
}

% --- Espacamento -----------------------------------------------------
\setlength{\parindent}{1.3cm}     % recuo de paragrafo
\setlength{\parskip}{0.2cm}       % espaco entre paragrafos
\OnehalfSpacing                   % entrelinhas 1,5, como pede a ABNT

% --- Comandos de conveniencia ----------------------------------------
% A abntex2 ja fornece \fonte{...} para a fonte de figuras e tabelas,
% exigida pela NBR 14724 logo abaixo do elemento. Use assim:
%   \fonte{Elaborado pelo autor (2026).}
% Adicione aqui os seus proprios comandos.
